\documentclass[10pt,a4paper]{amsart}
\usepackage[margin=19mm]{geometry}
\usepackage{amsmath,amssymb,mathtools}
\usepackage[scr=rsfs,cal=euler]{mathalfa}
\usepackage{xfrac}
\usepackage[unicode,hidelinks,bookmarks=false]{hyperref}
\newcommand{\ie}{i.e.\ }
\newcommand{\cf}{cf.\ }
\newcommand{\resp}{respectively\ }
\newcommand{\Subsection}{\subsection}
\providecommand{\nobreakdash}{\nobreak\hbox{-}}
\setlength{\emergencystretch}{2em}
\multlinegap0pt
\usepackage{amsthm}
\usepackage{braket}
\usepackage{dsfont}
\usepackage{mathtools}
\usepackage{enumitem}
\newtheorem{theorem}{Theorem}
\newtheorem{proposition}{Proposition}
\newtheorem{lemma}{Lemma}
\newtheorem{corollary}{Corollary}
\newtheorem{alemma}{Lemma}
\renewcommand{\thealemma}{B\arabic{alemma}}
\theoremstyle{definition}
\newtheorem{definition}{Definition}
\theoremstyle{remark}
\newtheorem{remark}{Remark}
\newcommand{\R}{\mathbb{R}}
\newcommand{\C}{\mathbb{C}}
\newcommand{\Z}{\mathbb{Z}}
\newcommand{\N}{\mathbb{N}}
\newcommand{\F}{\mathcal{F}}
\newcommand{\Hil}{\mathcal{H}}
\newcommand{\Op}{\mathcal{O}}
\newcommand{\K}{\mathcal{K}}
\newcommand{\one}{\mathds{1}}
\newcommand{\rmd}{\mathrm{d}}
\newcommand{\rmi}{\mathrm{i}}
\newcommand{\rme}{\mathrm{e}}
\newcommand{\bv}{\mathbf{v}}
\DeclareMathOperator{\supp}{supp}
\title[Newton-Cartan limit of Klein-Gordon AQFT]{Newton-Cartan limit of Klein-Gordon AQFT: spectral concentration and metastability of gravitational atoms (Corollary 1)}
\author[Leonardo Pachon]{Leonardo
Pachón}
\date{Source: arXiv:2604.26287v2 (CC BY 4.0)}
\begin{document}
\maketitle
\vspace{1em}
\noindent\textit{Source and licence.} Newton-Cartan limit of Klein-Gordon AQFT: gravitational atoms and the loss of vacuum entanglement. Version arXiv:2604.26287v2 (CC BY 4.0), distributed under the Creative Commons Attribution 4.0 International licence, http://creativecommons.org/licenses/by/4.0/, as recorded in the arXiv OAI licence metadata for this exact version and shown on the abstract page https://arxiv.org/abs/2604.26287v2. This item is an editorial re-presentation of the paper's spectral concentration and metastability corollary together with the paper's own proof of it and the source-local chain that proof uses. A full rights notice is delivered at licenses/arxiv-2604.26287-CC-BY-4.0.txt.
\noindent\textit{Editorial scope.} Counted unit: the paper's spectral concentration and metastability corollary for gravitational atoms, printed as Corollary 1 (source0.tex lines 1171--1196), delivered together with the paper's complete original proof of it (source0.tex lines 1198--1218) as the single primary proof body. No mathematics is written, repaired or re-derived: statement and proof are the paper's own text line for line, in the paper's own notation ($\mathcal{U}_c$, $P_c$, $H_S$, $E_n$, $F_c$, $r_{\rm s}$, $\alpha$, $T_{\rm B}$, $\Sigma_c$). Required context retained uncounted, because the counted proof is the paper's own assembly over it: the rescaled two-point reduction display (source0.tex lines 597--600); the limiting Schrodinger operator with its spectrum (lines 1060--1064); the statement of the Schwarzschild exterior theorem (lines 1071--1090); the tortoise-coordinate form of the kinetic operator (lines 2519--2522); and the approximants lemma with its complete proof (lines 2591--2655). Every cross-reference inside the retained spans resolves inside the delivered document, and no citation command occurs in any retained span, so the delivered document carries no bibliography and no reference or citation command of the delivered text was rewritten or adapted. Printed numbers are the paper's own: the wrapper restores the paper's per-kind theorem counters, its appendix lemma counter and its global equation counter before the retained material, so the retained Schwarzschild theorem prints as Theorem 2, the counted statement as Corollary 1, the retained approximants lemma as Lemma B.3, and the numbered displays of the retained spans carry the paper's own equation numbers (13), (23), (33) and (34). Omitted from this package and not redistributed: the paper's preamble and front matter as such (source0.tex lines 1--29 and 71--250 plus the closing material, of which only the title and the author name are reproduced in the wrapper); the introduction and the whole construction of the Newton-Cartan limit outside the retained ranges; the later theorems, propositions and corollaries; the appendices beyond the retained approximants lemma; and the paper's entire bibliography with its bib data file. Scope note, disclosed: the counted corollary is one of the paper's headline physical results and its written proof is short precisely because it is the top of a source-local chain; the chain it closes over is retained in full, namely the Schwarzschild exterior theorem whose parts (a) and (c) it invokes and the approximants lemma whose argument it repeats for the eigenfunction, so the counted proof can be checked line by line against the delivered text. Printed slips kept literal, among them the paper's own wording and its own choice of symbols. Layout and class: the paper's own class is the publisher's REVTeX class with AIP journal options, which is neither shipped with this package nor used to build it; the delivered wrapper is the lane's pinned amsart editorial wrapper at 10pt on A4, to which this package adds the paper's own theorem-like declarations with their separate per-kind counters, its own appendix lemma numbering and its own shorthand macros, all copied verbatim from the paper's source. No publisher class file, no publisher style file and no upstream script is redistributed. Assets: no retained span of this package contains a figure environment or a graphics-inclusion command, no image or artwork file exists in or is redistributed with this package, and the manifest and the delivery evidence both carry an empty asset-dependency list. A full rights notice is delivered at licenses/arxiv-2604.26287-CC-BY-4.0.txt.
\setcounter{equation}{12}
\begin{equation}\label{eq:H2}
\mathcal{U}_cF_f^{(c)}\longrightarrow F_f^{(\infty)}\quad\text{in }
L^2(\R^3),
\end{equation}

\setcounter{equation}{22}
\setcounter{theorem}{1}
\begin{equation}\label{eq:HS-schw}
H_S:=-\frac{\hbar^2}{2m}\Delta-\frac{GMm}{r}\quad\text{on }H^2(\R^3),
\qquad \sigma(H_S)=\Bigl\{E_n=-\frac{G^2M^2m^3}{2\hbar^2n^2}\Bigr\}_{n\geq1}
\cup[0,\infty).
\end{equation}

\begin{theorem}[Schwarzschild exterior]\label{thm:schw}
For the Schwarzschild exterior and $H_S$ as in~\eqref{eq:HS-schw}:
\begin{enumerate}[label=(\alph*)]
\item $\mathcal{U}_c(P_c-z)^{-1}\mathcal{U}_c^*\to(H_S-z)^{-1}$ strongly,
for every $z\in\C\setminus\R$. Hence $\mathcal{U}_c\,f(P_c)\,
\mathcal{U}_c^*\to f(H_S)$ strongly for every bounded continuous $f$,
and $\mathcal{U}_c\rme^{-\rmi tP_c/\hbar}\mathcal{U}_c^*\to
\rme^{-\rmi tH_S/\hbar}$ strongly, uniformly for $t$ in compact sets.
\item For every $c$, $\sigma(P_c)=[-mc^2,\infty)$ and $P_c$ has no
eigenvalues. The convergence in (a) is not in norm.
\item (Spectral concentration.) If $I$ is a bounded open interval with
$E_n\in I$ and $\bar I\cap\sigma(H_S)=\{E_n\}$, then $\mathcal{U}_c\one_I(P_c)
\mathcal{U}_c^*\to\Pi_n$ strongly, where $\Pi_n$ is the eigenprojection
of $H_S$ for $E_n$.
\item For $f\in C_c^\infty(\R\times(\R^3\setminus\{0\}))$ and $c$ large,
$F_f^{(c)}$ is defined, \eqref{eq:H2} holds, and the two-point functions
of $\hat\psi_c$ in the Boulware state converge to those of $\hat\psi$ in
the Fock vacuum of $\rmd\Gamma(H_S)$.
\end{enumerate}
\end{theorem}

\begin{corollary}[Spectral concentration and metastability]
\label{cor:metastability}
Let $u$ be a normalised eigenfunction of $H_S$ for $E_n$. Then
\[
\braket{\mathcal{U}_c^*u,\rme^{-\rmi tP_c/\hbar}\mathcal{U}_c^*u}
\longrightarrow\rme^{-\rmi tE_n/\hbar}
\]
uniformly for $t$ in compact sets, and the spectral measure of $\mathcal{U}_c^*u$ for $P_c$ concentrates at $E_n$:
its mass in every open interval containing $E_n$ tends to $1$. More
precisely, with $\alpha=GMm/\hbar c$ and the gravitational Bohr time
$T_{\rm B}:=\hbar^3/G^2M^2m^3$, there is $C$ such that for $c$ large and
all $t$
\[
\bigl|\braket{\mathcal{U}_c^*u,\rme^{-\rmi tP_c/\hbar}\mathcal{U}_c^*u}
-\rme^{-\rmi tF_c(E_n)/\hbar}\bigr|\leq C\bigl(1+|t|/T_{\rm B}\bigr)\,
\alpha^{1/2},
\]
and the spectral measure of $\mathcal{U}_c^*u$ for $P_c$ gives mass at
most $C\alpha(mc^2\alpha^2/\delta)^2+C\alpha^2$ to the complement of
$[F_c(E_n)-\delta,F_c(E_n)+\delta]$.
For every fixed $c$, on the other hand,
\[
\lim_{T\to\infty}\frac1T\int_0^T\bigl|\braket{\mathcal{U}_c^*u,
\rme^{-\rmi tP_c/\hbar}\mathcal{U}_c^*u}\bigr|^2\rmd t=0 .
\]
\end{corollary}

\begin{proof}
The first two statements are Theorem~\ref{thm:schw}(a) and (c). For the
bounds, the proof of Lemma~\ref{lem:approximants} applies to
$\varphi=u$, for which $|u|+|\partial_ru|+|\partial_r^2u|+|\Delta u|\leq
C(1+r^{-1})\rme^{-r/2na_0}$, with $u_c:=\theta(r_*/L)\,u\in D(S_c)$ in
place of an element of $C_c^\infty(\Sigma_c)$; it gives
$\|\mathcal{U}_cu_c-u\|=O(r_{\rm s})$ and
$\|(S_c-E_n)u_c\|=\|\mathcal{U}_cS_cu_c-E_n\mathcal{U}_cu_c\|=
O(r_{\rm s}^{1/4})=O(c^{-1/2})$. Since
$|F_c(s)-F_c(E_n)|\leq2|s-E_n|/(1+2E_n/mc^2)^{1/2}$ for $s\geq-mc^2/2$,
also $\|(P_c-F_c(E_n))u_c\|=O(c^{-1/2})$, and Duhamel's formula gives
$|\braket{u_c,\rme^{-\rmi tP_c/\hbar}u_c}-\rme^{-\rmi tF_c(E_n)/\hbar}
\|u_c\|^2|\leq|t|\,\|(P_c-F_c(E_n))u_c\|\,\|u_c\|/\hbar$. Finally
$\|u_c-\mathcal{U}_c^*u\|\leq\|\mathcal{U}_cu_c-u\|$ and
$\|u_c\|^2=1+O(r_{\rm s})$. The spectral mass of $u_c$ outside the
interval is at most $\|(P_c-F_c(E_n))u_c\|^2/\delta^2$, and that of
$\mathcal{U}_c^*u$ differs from it by $O(r_{\rm s})$. For $l\geq1$ the
approximants do better, $u\sim r^l$ making the Newtonian term near the
horizon smaller, and the rate is $\alpha$. The last statement is Wiener's
theorem, since $P_c$ has no eigenvalues.
\end{proof}

\setcounter{equation}{32}
\setcounter{alemma}{2}
\begin{equation}\label{eq:Ac-tortoise}
A_c=-\frac{1}{r^2}\,\partial_*\bigl(r^2\partial_*\,\cdot\,\bigr)-
\frac{N_c^2}{r^2}\,\Delta_{S^2}.
\end{equation}

\begin{alemma}[Approximants]\label{lem:approximants}
Let $\theta\in C^\infty(\R)$ with $0\leq\theta\leq1$, $\theta=0$ on
$(-\infty,-2]$ and $\theta=1$ on $[-1,\infty)$, let $L=(\ell r_{\rm s})^{1/2}$ for a fixed
length $\ell$, and for $\varphi\in C_c^\infty(\R^3)$ put $u_c:=\theta(r_*/L)\,
\varphi|_{\Sigma_c}$. Then $u_c\in C_c^\infty(\Sigma_c)$,
$\|\mathcal{U}_cu_c-\varphi\|=O(r_{\rm s})$ and $\|\mathcal{U}_cS_cu_c-
H_S\varphi\|=O(r_{\rm s}^{1/4})$, where $\mathcal{U}_cS_cu_c$ is extended
by zero to $\{r\leq r_{\rm s}\}$.
\end{alemma}

\begin{proof}
Measure lengths in units of $\ell$, so that $L=r_{\rm s}^{1/2}$. Let $C$
denote constants depending on $\varphi$, $\theta$ and $\ell$ only, and
let $r_{\rm s}\leq\frac14$. Two facts about $r_*$ are used. First,
$r_*(2r_{\rm s})=2r_{\rm s}>-L$, so $\theta(r_*/L)=1$ for $r\geq
2r_{\rm s}$. Second, where $r_*\leq-L$ one has $\ln((r-r_{\rm s})/r_{\rm s})
\leq-L/r_{\rm s}-1$, so on the support of $\theta'(r_*/L)$
\begin{equation}\label{eq:tiny-lapse}
N_c^2\leq\rme^{-1-r_{\rm s}^{-1/2}} .
\end{equation}
The region $\{-2L\leq r_*\leq2r_{\rm s}\}$, which contains the part of
$\supp u_c$ with $r<2r_{\rm s}$, has tortoise length at most $3L$.

\emph{The vectors.} $\mathcal{U}_cu_c=N_c^{-1}\theta\varphi$ on
$\Sigma_c$. On $\{r\leq r_{\rm s}\}$, $\|\varphi\|^2\leq Cr_{\rm s}^3$. On
$\{r\geq2r_{\rm s}\}$, $|N_c^{-1}-1|\leq\sqrt2\,r_{\rm s}/r$, contributing
$Cr_{\rm s}^2$ to the square of the norm. On $\{r_{\rm s}<r<2r_{\rm s}\}$
we use $\rmd^3x=N_c^2\,\rmd\mu_c$: $\int N_c^{-2}\theta^2|\varphi|^2
\rmd^3x=\int\theta^2|\varphi|^2r^2\,\rmd r_*\rmd\Omega\leq Cr_{\rm s}^2L$,
and $\int|\varphi|^2\rmd^3x\leq Cr_{\rm s}^3$.

\emph{The operator.} Write $S_cu_c=\theta S_c\varphi+\frac{\hbar^2}{2m}
[A_c,\theta]\varphi$. For smooth $\varphi$, with $\Delta_r=\partial_r^2+
\frac2r\partial_r$ and $\Delta=\Delta_r+r^{-2}\Delta_{S^2}$,
\[
N_c^{-1}S_c\varphi-H_S\varphi=-\frac{\hbar^2}{2m}\Bigl[(N_c-1)\Delta
\varphi-N_c\frac{r_{\rm s}}{r}\Delta_r\varphi+N_c\frac{r_{\rm s}}{r^2}
\partial_r\varphi\Bigr]-\frac{GMm}{r}\bigl(N_c^{-1}-1\bigr)\varphi .
\]
Here $|\Delta\varphi|\leq C$, $|\Delta_r\varphi|+|\partial_r\varphi|
\leq C(1+r^{-1})$, $0\leq1-N_c\leq r_{\rm s}/r$ and $N_c\leq1$. The
first three terms have $L^2$ norm $O(r_{\rm s}^{1/2})$ on $\Sigma_c$. The
last is $O(r_{\rm s}^{1/2})$ on $\{r\geq2r_{\rm s}\}$, and on
$\{r_{\rm s}<r<2r_{\rm s}\}$ its square norm, weighted by $\theta^2$, is
at most $C\int\theta^2N_c^{-2}\,\rmd^3x/r^2\leq C\int\theta^2\,\rmd r_*
\leq CL$. Moreover $(\theta-1)H_S\varphi$ and $H_S\varphi$ restricted to
$\{r\leq r_{\rm s}\}$ are supported in $\{r<2r_{\rm s}\}$, where
$\|H_S\varphi\|^2\leq Cr_{\rm s}$. For the commutator,
by~\eqref{eq:Ac-tortoise},
\[
[A_c,\theta]\varphi=-\frac{1}{L^2}\theta''\varphi-\frac2L\theta'\,
\partial_*\varphi-\frac{2N_c^2}{rL}\theta'\varphi,
\]
with $\theta$ and its derivatives evaluated at $r_*/L$. It is supported
in $\{-2L\leq r_*\leq-L\}$, where $r<2r_{\rm s}$ and
\eqref{eq:tiny-lapse} holds, and $\partial_*\varphi=N_c^2\partial_r
\varphi$. Its norm in $\K_c$, which is the norm of $\mathcal{U}_c$
applied to it, satisfies
\[
\bigl\|[A_c,\theta]\varphi\bigr\|^2\leq C\,r_{\rm s}^2\int_{-2L}^{-L}
\Bigl(\frac{1}{L^2}+\frac{\rme^{-r_{\rm s}^{-1/2}}}{r_{\rm s}L}\Bigr)^2
\rmd r_*\leq C\frac{r_{\rm s}^2}{L^3}=C\,r_{\rm s}^{1/2}.
\]
Collecting terms, $\|\mathcal{U}_cS_cu_c-H_S\varphi\|=O(r_{\rm s}^{1/4})$.
\end{proof}

\end{document}
